% !TEX program = lualatex
\DocumentMetadata{pdfstandard=UA-2,pdfversion=2.0,lang=en-US,tagging=on}
\documentclass{./../../latex/tests}
\begin{document}
\thispagestyle{plain}
\myheader{Midterm Sample: Solutions}
\rhead{Midterm Sample: Solutions}
\vspace{0.5em}

\begin{enumerate}

%%%%%%%%%%%%%% Question 1
\item (5 pts)
\begin{enumerate}
\item (1 pt) \textit{Is $f$ a valid function?} \textbf{Yes}: a function may map two values of $x$ to the same value of $y$; what it cannot do is map one $x$ to two values of $y$.
\item (1 pt) $A \cup B =$ Set of all integers.
\item (1 pt) $A \cup B = A \quad \quad A \cap B = B$
\item (1 pt) No: since the function is not one-to-one (it is not monotonic), the inverse does not exist.
\item (1 pt) $(0, 2)$
\end{enumerate}

%%%%%%%%%%%%%% Question 2
\item (5 pts) $A = I - X(X^T X)^{-1}X^T$
\begin{enumerate}
\item (3 pts) Say the dimension of $X$ is \( m \times n \). Then the dimension of \( X^T_{n \times m} X_{m \times n} \) is \( n \times n \), so the dimension of \( (X^T X)^{-1} \) is also \( n \times n \). This implies that the dimension of \( X_{m \times n}(X^T X)^{-1}_{n \times n} X^T_{n \times m} \) is \( m \times m \). Hence $A$ is square (as is $X^T X$), but $X$ need not be square.
\item (2 pts)
\[
\begin{aligned}
A A &= (I-X(X^T X)^{-1} X^T)(I-X(X^T X)^{-1} X^T) \\
& = I-X(X^T X)^{-1} X^T-X(X^T X)^{-1} X^T+X\underbrace{(X^T X)^{-1} X^T X}_{I}(X^T X)^{-1} X^T \\
&= I-X(X^T X)^{-1} X^T = A
\end{aligned}
\]
\end{enumerate}

%%%%%%%%%%%%%% Question 3
\item (10 pts)
\begin{enumerate}
\item (2 pts) Setting supply equal to demand in each market:
\[ 30 + p_e = 100 - 2p_e + p_g \implies 3p_e - p_g = 70 \]
\[ 20 + p_g = 50 + p_e - p_g \implies -p_e + 2p_g = 30 \]
\item (2 pts)
\[ A = \begin{bmatrix} 3 & -1 \\ -1 & 2 \end{bmatrix}_{2 \times 2} \quad
x = \begin{bmatrix} p_e \\ p_g \end{bmatrix}_{2 \times 1} \quad
b = \begin{bmatrix} 70 \\ 30 \end{bmatrix}_{2 \times 1} \]
\item (2 pts) The cofactors are $|C_{11}| = 2$, $|C_{12}| = 1$, $|C_{21}| = 1$, $|C_{22}| = 3$, so
\[ adj A = C^T = \begin{bmatrix} 2 & 1 \\ 1 & 3 \end{bmatrix} \]
\item (1 pt) $|A| = 3 \times 2 - (-1) \times (-1) = 5 \neq 0$, so $A$ is nonsingular.
\item (1 pt) $A^{-1} A x = A^{-1} b \implies x^* = A^{-1} b$
\item (2 pts)
\[ x^* = A^{-1} b = \frac{1}{|A|} adj A \, b
= \frac{1}{5}\begin{bmatrix} 2 & 1 \\ 1 & 3 \end{bmatrix}\begin{bmatrix} 70 \\ 30 \end{bmatrix}
= \frac{1}{5}\begin{bmatrix} 170 \\ 160 \end{bmatrix}
= \begin{bmatrix} 34 \\ 32 \end{bmatrix} \]
So $p_e^* = 34$ and $p_g^* = 32$ (and, substituting back, $q_e^* = 64$, $q_g^* = 52$).
\end{enumerate}

\end{enumerate}

\end{document}
